\problemname{Nice Numbers}
The Swedish junior programming team loves having a nice time when they
are out traveling. One who appreciated niceness more than anyone else was Mårten.
However, he did not share the same view of niceness as everyone else. For Mårten,
niceness had a purely mathematical definition: a number is nice if and
only if it is divisible by 3.

Since not all numbers are nice, one must of course find ways to make them
nice. You will be given an integer, and your task is to help Mårten
figure out in how many ways he can make the number nice by deleting certain
of its digits. He may not delete all the digits, and the resulting
number may not have leading zeros (however, the number $0$ is considered nice,
so a single zero is fine). Remember that a number is divisible by three if
and only if its digit sum is divisible by three.

\section*{Input}
The first and only line of input contains an integer (with up to $100\,000$
digits). The number contains only digits $0$--$9$. The number does not start with a zero.

\section*{Output}
Your program should output a single number on a line: the number of different ways
Mårten can remove digits from the number so that it becomes divisible by $3$. Two ways
are considered different if the sets of positions where digits are removed differ. Since
the answer can be very large, output the remainder when the answer is divided by one billion ($10^9$).

\section*{Scoring}
Your solution will be tested on a set of test groups.
To earn points for a group, you must pass all test cases in that group.

\noindent
\begin{tabular}{| l | l | l |}
  \hline
  \textbf{Group} & \textbf{Points} & \textbf{Constraints} \\ \hline
  $1$   & $20$       & $N \leq 6$ and the answer is less than one billion \\ \hline
  $2$   & $20$       & $N \leq 20$ and the answer is less than one billion \\ \hline
  $3$   & $20$       & $N \leq 1000$ \\ \hline
  $4$   & $40$       & No additional constraints. \\ \hline
\end{tabular}

\section*{Explanation of Samples}
Explanation of sample 1: Mårten can achieve divisibility in one way, by deleting the one. What remains is just the number 3, which of course is divisible by three.

Explanation of sample 2: Mårten achieves divisibility by not deleting any digit at all.

Explanation of sample 3: it is not possible to create a nice number regardless of what Mårten does.

Explanation of sample 4: one can create the nice numbers 9, 192, 912, and 12, but the last one can be achieved in two ways: by deleting the first two digits or by deleting the middle two digits.

Explanation of sample 5: one can create the nice numbers 0, 3, 12, 102, 123, and 1023.
